\section*{अध्याय \ref{ch:b1:first-law} --- ऊष्मागतिकी का पहला नियम}

\begin{solution}{exo:b1:first-law:1}
$Q = C_{P,m}\Delta T = 29.1 \times 100 = \qty{2.91}{kJ}$; $W = -nR\Delta T =
\qty{-0.83}{kJ}$; $\Delta U = Q + W = \qty{2.08}{kJ} = C_{V,m}\Delta T$।
\end{solution}

\begin{solution}{exo:b1:first-law:2}
$W = -nRT\ln(V_2/V_1) = -2494\ln(0.25) = \qty{+3.46}{kJ}$; $\Delta U = 0$, अतः
$Q = \qty{-3.46}{kJ}$ (जो \omterm{def:b1:first-law:heat}{ऊष्मा-भंडार} को दी गई)।
\end{solution}

\begin{solution}{exo:b1:first-law:3}
$T_2 = T_1 \times 10^{\gamma - 1} = 300 \times 10^{0.4} = \qty{754}{K}$; $P_2 =
10^{1.4} = \qty{25}{bar}$।
\end{solution}

\begin{solution}{exo:b1:first-law:4}
आर्गन: $C_{V,m} = \tfrac32 R = 12.5$, $C_{P,m} = \tfrac52 R = 20.8$ जूल प्रति मोल-केल्विन,
$\gamma = 1.67$; नाइट्रोजन: $20.8$, $29.1$, $\gamma = 1.40$। पानी के लिए गरम
होने पर प्रसार बहुत छोटा है, अतः स्थिर दाब पर कार्य $P\,\dd V$ \omterm{def:b1:first-law:heat}{ऊष्मा} के
सामने नगण्य रहता है: $c_P - c_V \approx 0$।
\end{solution}

\begin{solution}{exo:b1:first-law:5}
$T_f = (200 \times 80 + 300 \times 20)/500 = \qty{44}{\degreeCelsius}$। प्याले
के साथ: $(200 \times 4.18 \times 80 + (300 \times 4.18 + 80) \times 20)/(500 \times
4.18 + 80) = \qty{43}{\degreeCelsius}$।
\end{solution}

\begin{solution}{exo:b1:first-law:6}
बर्फ़: $0$ तक गरम करना: $0.050 \times 2100 \times 10 = \qty{1.05}{kJ}$; गलाना:
\qty{16.7}{kJ}; कुल \qty{17.8}{kJ}। पानी $0$ तक ठंडा होकर \qty{25.1}{kJ} दे
सकता है: अतः सारी बर्फ़ गल जाती है; फिर $0.250 \times 4180\,T_f = 25.1 - 17.8$ kJ: $T_f =
\qty{7}{\degreeCelsius}$। \qty{150}{g} के साथ: \qty{53.3}{kJ} चाहिए, $>$ \qty{25.1}{kJ}:
अतः केवल $(25.1 - 3.15)/334 = \qty{66}{g}$ बर्फ़ गलती है, $T_f = \qty{0}{\degreeCelsius}$,
और \qty{84}{g} बर्फ़ बची रह जाती है।
\end{solution}

\begin{solution}{exo:b1:first-law:7}
$W = -P_1(V_2 - V_1) + 0 - P_2(V_1 - V_2) + 0 = (P_2 - P_1)(V_2 - V_1) > 0$:
चक्र घड़ी की उलटी दिशा में तय होता है (कम दाब पर प्रसार, ऊँचे दाब पर
संपीडन), अतः गैस को परिणामी कार्य \emph{मिलता} है, जो घिरे हुए क्षेत्रफल के
बराबर है; घड़ी की दिशा में वह उसे देती --- यानी एक इंजन।
\end{solution}

\begin{solution}{exo:b1:first-law:8}
$T_2 = 300 \times 7^{0.286} = \qty{523}{K}$। गरम संपीडित हवा हर झटके पर चालन
से नली को गरम कर देती है। टायर में हवा स्थिर आयतन पर परिवेश के ताप तक ठंडी
होती है: $P \propto T$, अतः दाब तापों के अनुपात में गिर जाता है --- उसे फिर
भर लीजिए।
\end{solution}

\begin{solution}{exo:b1:first-law:9}
$W = 0$ (धकेलने के लिए कोई बाहरी दाब ही नहीं), $Q = 0$: अतः $\Delta U = 0$,
और \omterm{def:b1:kinetic-theory:model}{आदर्श गैस} के लिए $\Delta T = 0$। अनुत्क्रमणीय: गैस अपने आप कभी वापस नहीं
बहती। कोई वास्तविक गैस, जिसके अणु आपस में आकर्षित होते हैं, थोड़ी ठंडी हो
जाती है (अलग होते समय \omterm{thm:b1:work-and-energy:ke}{गतिज ऊर्जा} का कुछ भाग \omterm{def:b1:work-and-energy:conservative}{स्थितिज ऊर्जा} बन जाता है)।
\end{solution}

\begin{solution}{exo:b1:first-law:10}
$Q = L_v = \qty{2.26}{MJ}$; $W = -P(V_g - V_l) = -1.013\times10^5 \times 1.67 =
\qty{-0.17}{MJ}$; $\Delta U = \qty{2.09}{MJ}$: यानी \omterm{def:b1:first-law:heat}{ऊष्मा} का $93\%$ अणुओं को
अलग करने में जाता है (\omterm{def:b1:kinetic-theory:U}{आंतरिक ऊर्जा}), और $7\%$ वायुमंडल को पीछे धकेलने में।
\end{solution}

\begin{solution}{exo:b1:first-law:11}
$W = -P_2(V_2 - V_1)$, $\Delta U = nC_{V,m}(T_2 - T_1)$, $V_2 = nRT_2/P_2$, $V_1
= nRT_1/P_1$: $\tfrac52(T_2 - 300) = -T_2 + 300 \times 5$, $3.5T_2 = 2250$,
$T_2 = \qty{643}{K}$, $V_2 = \qty{10.7}{L}$। उत्क्रमणीय: $T_2 = 300 \times
5^{0.286} = \qty{475}{K}$, $V_2 = \qty{7.9}{L}$। अचानक किया गया संपीडन गैस पर
अधिक कार्य करता है (पूरी \qty{5}{bar} शुरू से ही लगती है) और अंत में गैस
अधिक गरम तथा कम संपीडित रहती है।
\end{solution}

\begin{solution}{exo:b1:first-law:12}
$\rho = PM/RT = \qty{1.20}{kg/m^3}$: $\sqrt{\gamma P/\rho} = \sqrt{1.4 \times
1.013\times10^5/1.20} = \qty{343}{m/s}$; $\sqrt{P/\rho} = \qty{290}{m/s}$, यानी $15\%$
कम। $P/\rho = RT/M$ के साथ: $c = \sqrt{\gamma RT/M}$; और $u = \sqrt{3RT/M}$
= \qty{502}{m/s} के सामने: $c/u = \sqrt{\gamma/3} = 0.68$।
\end{solution}

\begin{solution}{pb:b1:first-law:1}
\textbf{1.} $n = P_1V_1/RT_1 = 10^5 \times 2.12\times10^{-4}/2494 = \qty{8.5e-3}{mol}$।

\textbf{2.} $PV^\gamma$ अचर: $V_2 = V_1(P_1/P_2)^{1/\gamma} = 212 \times
7^{-0.714} = \qty{53}{cm^3}$।

\textbf{3.} $T_2 = T_1(P_2/P_1)^{(\gamma - 1)/\gamma} = 300 \times 7^{0.286} =
\qty{523}{K}$।

\textbf{4.} $W = nC_{V,m}(T_2 - T_1) = 8.5\times10^{-3} \times 20.8 \times 223 =
\qty{39}{J}$।

\textbf{5.} जोड़ी जाने वाली हवा: $\Delta n = \Delta P\,V/RT = 6\times10^5 \times 2\times10^{-3}/
2494 = \qty{0.48}{mol}$: यानी $0.48/8.5\times10^{-3} = 57$ झटके, लगभग
\qty{2.2}{kJ} --- एक मिनट की ईमानदार मेहनत।

\textbf{6.} स्थिर आयतन पर $P \propto T$: हवा के ठंडा होते ही दाब गिर जाता
है (अंतिम झटके की हवा के लिए $300/523$ तक); अतः साइकिल-सवार कुछ देर रुककर
कुछ और झटके लगाता है।

\textbf{7.} $W_{\text{समतापी}} = nRT\ln(P_2/P_1) = 8.5\times10^{-3} \times 2494 \times
1.95 = \qty{41}{J} > \qty{39}{J}$: समतापी पथ \qty{7}{bar} तक पहुँचने के लिए
गैस को और अधिक संपीडित करता है (\qty{53}{cm^3} की जगह \qty{30}{cm^3} तक),
क्योंकि वह ठंडी बनी रहती है --- अधिक आयतन बुहारा गया, अधिक कार्य।

\textbf{8.} $n = 10^5 \times 5\times10^{-4}/2494 = \qty{0.020}{mol}$; $T_2 =
300 \times 20^{0.4} = \qty{994}{K}$।

\textbf{9.} $P_2 = 20^{1.4} = \qty{66}{bar}$।

\textbf{10.} \qty{994}{K} ईंधन के स्वतः प्रज्वलन बिंदु से कहीं ऊपर है: अतः
अंतःक्षिप्त ईंधन संस्पर्श होते ही जल उठता है। कोई पेट्रोल इंजन अपने
हवा--ईंधन मिश्रण को केवल लगभग \qty{750}{K} तक ही संपीडित करता है, ताकि वह
स्फुलिंग से पहले न जले; अनुपात बहुत ऊँचा हो तो ``खटखट'' होने लगती है।

\textbf{11.} $W = nC_{V,m}(T_2 - T_1) = 0.020 \times 20.8 \times 694 = \qty{289}{J}$।

\textbf{12.} प्रति बेलन $25 \times 289 = \qty{7.2}{kW}$, जो गरम गैस में संचित
रहती है और प्रसार में लौटा दी जाती है।

\textbf{13.} अनुत्क्रमणीयता (गैस अर्ध-स्थैतिक पथ पर नहीं चलती; अतिरिक्त
कार्य क्षय हो जाता है) $T_2$ को बढ़ा देती है; और दीवारों को खोई \omterm{def:b1:first-law:heat}{ऊष्मा} उसे
घटा देती है। किसी वास्तविक इंजन में दूसरा प्रायः जीत जाता है: $T_2$ आदर्श
मान से कुछ नीचे रहता है।

\textbf{14.} $Q = 20\times10^{-6} \times 43\times10^6 = \qty{860}{J}$।

\textbf{15.} समदाबी: $Q = nC_{P,m}(T_3 - T_2)$: $T_3 - T_2 = 860/(0.020 \times
29.1) = \qty{1480}{K}$, $T_3 = \qty{2470}{K}$; $V_3 = V_2T_3/T_2 = 25 \times
2.48 = \qty{62}{cm^3}$।

\textbf{16.} $W = -P_2(V_3 - V_2) = -66\times10^5 \times 37\times10^{-6} =
\qty{-244}{J}$: यानी गैस द्वारा दिया गया कार्य।

\textbf{17.} $T_4 = T_3(V_3/V_1)^{\gamma - 1} = 2470 \times (0.124)^{0.4} =
\qty{1070}{K}$; $W = nC_{V,m}(T_4 - T_3) = 0.020 \times 20.8 \times (-1400) =
\qty{-582}{J}$।

\textbf{18.} दिया गया परिणामी कार्य $= 244 + 582 - 289 = \qty{537}{J}$; अनुपात
$537/860 = 62\%$ (इन्हीं मानों पर आदर्श डीज़ल दक्षता; वास्तविक इंजन लगभग
$40\%$ तक पहुँचते हैं)।

\textbf{19.} $1 - 300/2470 = 88\%$: आदर्श डीज़ल चक्र कार्नो सीमा से कहीं
नीचे रह जाता है, क्योंकि उसे \omterm{def:b1:first-law:heat}{ऊष्मा} तापों की एक परास पर मिलती है, न कि
सबसे ऊँचे ताप पर।

\textbf{20.} किसी चक्र पर $\Delta U = 0$: $Q_{\text{बाहर}} = -(860 - 537) =
\qty{-323}{J}$, अर्थात् \qty{323}{J} निकास के साथ चली जाती है; जाँच:
$nC_{V,m}(T_4 - T_1) = 0.020 \times 20.8 \times 770 = \qty{320}{J}$।

\textbf{21.} \qty{10}{L} हवा में ऊष्मा-चालन की तुलना में दोलन तेज़ है (एक
सेकंड): अतः रुद्धोष्म। $PV^\gamma$ को रैखिक करने पर:
$\dd P/P = -\gamma\,\dd V/V = -\gamma Sx/V$।

\textbf{22.} गोली पर बल $S\,\dd P = -\gamma PS^2x/V$: $m\ddot x =
-(\gamma PS^2/V)x$, यानी सरल आवर्त, $\omega_0^2 = \gamma PS^2/mV$।

\textbf{23.} $\omega_0^2 = 1.4 \times 10^5 \times 4\times10^{-8}/(0.0165 \times 0.010)
= \qty{33.9}{s^{-2}}$: $T = 2\pi/5.83 = \qty{1.08}{s}$। $T = \qty{1.10}{s}$ से:
$\gamma = 4\pi^2mV/(T^2PS^2) = 4\pi^2 \times 0.0165 \times 0.010/(1.21 \times 10^5
\times 4\times10^{-8}) = 1.35$।

\textbf{24.} घर्षण गति को अवमंदित करता है और, यदि वह सूखा हो, तो साम्य को
भी खिसका देता है; और गोली के पास से रिसाव संपीडन के दौरान हवा निकल जाने
देता है, जिससे प्रत्यानयन बल घट जाता है: दोनों आवर्तकाल को लंबा करते हैं
और $\gamma$ को नीचे की ओर झुका देते हैं --- जैसे ऊपर का $1.35$।

\textbf{25.} पंप और डीज़ल: $Q = 0$, अतः $\Delta U = W$ --- यानी कार्य \omterm{def:b1:kinetic-theory:U}{आंतरिक
ऊर्जा} और ताप बन गया, चरघातांक $\gamma$ वाले लाप्लास के नियम के अनुसार।
गोली: कोई तेज़ संपीडन रुद्धोष्म होता है, और $\gamma$ किसी वायु-स्प्रिंग की
दृढ़ता तय करता है।
\end{solution}
